Existing benchmarks evaluate code generation, issue resolution, static method reachability, and runtime call-structure recovery. In this work, we propose \benchmark, a dataset of 1,232 generated Java classes and 46,656 questions about whether a target method executes from a source under values supplied by a generated execution. A separate configurable algorithm constructs this dataset by varying class size, guarded calls, and values passed between methods. This allows us to compare model behavior across structural difficulties, with labels and executable graphs validated by running the complete dataset. The model returns a Boolean answer and a method trace, which are evaluated separately against the label and executable calls. A positive trace must reach the target, while a negative trace must cover all methods reachable from the source. On the full Class-only (CO) dataset, Qwen~3.6 obtains 87.98\% answer correctness and 67.43\% trace correctness, compared with 84.21\% and 55.90\% for Gemma~4. A correct answer therefore does not always come with a valid trace, particularly when questions involve guarded calls, passed values, or complete exploration for negative targets. The Explicit execution (EE) variant uses the same dataset but makes the entry point and deterministic execution explicit. Qwen~3.6 obtains slightly higher answer and trace scores with EE, although its reasoning budget also increases. On 216 difficult questions evaluated at the same budget, Qwen3.8-Flash-Next's trace correctness increases from 35.65\% to 46.76\%, while non-responses increase from 10.65\% to 19.44\%.
